\documentclass[a4paper,11pt]{article}

\usepackage{amsmath, amssymb, amstext, amsfonts, mathrsfs}


% \sffamily %schrift ohne Serifen

\usepackage[T1]{fontenc} 
% schriftencodierung f�r umlaute, trennung
% f\"ur Uni
\usepackage[latin1]{inputenc}
\usepackage{selinput}
% \usepackage[utf8x]{inputenc} 
\usepackage{bibgerm} 
% german bibliography
\usepackage[german]{babel}
%wichtig f�r deutschen Content
\usepackage{ucs}
%erweiterte UTF-8 Unterst�tzung
\usepackage{wrapfig} 
% Paket zur Positionierung einbinden
\usepackage{multirow}
% zusammenfassen von Tabellenzellen
% \usepackage{subscript}
% zum tiefstellen
\usepackage{lscape}
\usepackage{pdflscape}
% zum drehen der Seite
% \usepackage[super]{natbib}
\usepackage[square,sort,comma,numbers]{natbib}
% Erstellung es Literaturverzeichnisses
\usepackage{url}
% Umbruch f�r URL
\usepackage{pst-3dplot}
% f�r tex Grafiken n�tig
\usepackage{pstricks}
\usepackage{listings}
% f�r das einf�gen von Quelltext
\definecolor{codegray}{rgb}{0.92,0.92,0.92}
\lstset{basicstyle=\fontsize{9}{11}\selectfont\ttfamily, breaklines=true, backgroundcolor=\color{codegray}, numbers=left, numberstyle=\tiny, tabsize=4, language=java}
\definecolor{mymauve}{rgb}{0.58,0,0.82}
\definecolor{mygreen}{rgb}{0,0.6,0}
\lstset{
commentstyle=\color{mygreen},
keywordstyle=\color{mymauve},
language=Java,
stringstyle=\color{blue}
}


%Einstellungen f�r Quellcode
\usepackage[a4paper, left=3cm, right=2cm, top=2cm]{geometry}
% Formatierung R�nder
\usepackage[section]{placeins}
% f�r Floatbarriere
\usepackage{color}
\usepackage{colortbl}
%f�r die Verwendung von Farben

\clubpenalty = 10000 
\widowpenalty = 10000
\displaywidowpenalty = 10000
%Verhinderung von Hurenkindern und Schusterjungen
%10000 bedeutet die sie sollen kommplett vermieden werden

\title{Entwicklung einer Android-Applikation f�r die Alarmierung der Einsatzkr�fte der Freiwilligen Feuerwehren}

\author{Sebastian Rieger}

\pagenumbering{arabic}
%Seitenzahlen(arabische Zahlen)

\setlength{\parindent}{0.25cm} 
%Absatzeinzug �ndern in Zoll
\setlength{\parskip}{0.25cm}
%Absatzabstand
\linespread {1.5}
%Zeilenabstand

\usepackage{setspace}
\usepackage{hyperref}
%anklickbare Hyperlinks

%funktioniert nicht bei Fu�noten
\usepackage{graphicx}
\usepackage{graphics}
%f�r einbinden von Grafiken

\usepackage{framed}
%f�r Umrandung der Erkl�rung
\usepackage{acronym}
% f�r abk�rzungen
% \usepackage{PSTricks}
\usepackage{epstopdf}
% f�r eps bilder nutze pdflatex --shell-escape this.tex
\usepackage{amssymb}
% f�r mathematische Symbole

\usepackage{hyperref}
% klickbare links

% \usepackage{pdfpages}
\usepackage{rotating}
\usepackage{svg}
%%%%%%%%%%%%%%%%%%%%%%%%%%%%%%%%%%%%%%%%%%%%%%%%%%%%%%%%%%%%%%%%%%%%%%%%%%%%%%%%%%%%%%%%%%%%%%%%%%%%%
%% Angaben zur Arbeit
%%%%%%%%%%%%%%%%%%%%%%%%%%%%%%%%%%%%%%%%%%%%%%%%%%%%%%%%%%%%%%%%%%%%%%%%%%%%%%%

\newcommand{\Autor}{Sebastian Rieger}
\newcommand{\MatrikelNummer}{10286908}
\newcommand{\Kursbezeichnung}{TINF12B1}

\newcommand{\FirmenName}{PDV Systeme}
\newcommand{\FirmenStadt}{Erfurt}
\newcommand{\FirmenLogoDeckblatt}{{\includegraphics[width=3cm]{}}}

% Falls es kein Firmenlogo gibt:
%  \newcommand{\FirmenLogoDeckblatt}{}

\newcommand{\BetreuerFirma}{Prof. Rolf Kruse}
\newcommand{\BetreuerDHBW}{Prof. Steffen Avemarg}
\newcommand{\Titel}{Entwicklung eines Natural User Interface mit Hilfe moderner AR und AI Technologien}
\newcommand{\AbgabeDatum}{15.11.2017}

\newcommand{\Dauer}{24 Wochen}

% \newcommand{\Abschluss}{Bachelor of Engineering}
\newcommand{\Abschluss}{Master of Science}

\newcommand{\Studiengang}{Angewandte Informatik}
% \newcommand{\Studiengang}{Angewandte Informatik}
\newcommand{\Was}{Masterarbeit}

%%%%%%%%%%%%%%%%%%%%%%%%%%%%%%%%%%%%%%%%%%%%%%%%%%%%%%%%%%%%%%%%%%%%%%%%%%%%%%%%%%%%%%%%%%%%%%%%%%%%% 
%steuervariable
\usepackage{ifthen} %Package f�r if/else
\newboolean{bilder} %Deklaration
\setboolean{bilder}{true} %Zuweisung
% \setboolean{bilder}{false} %Zuweisung
%%%%%%%%%%%%%%%%%%%%%%%%%%%%%%%%%%%%%%%%%%%%%%%%%%%%%%%%%%%%%%%%%%%%%%%%%%%%%%%%%%%%%%%%%%%%%%%%%%%%%

\lstdefinelanguage{JavaScript}{
  keywords={typeof, new, true, false, catch, function, return, null, catch, switch, var, if, in, while, do, else, case, break},
  keywordstyle=\color{blue}\bfseries,
  ndkeywords={class, export, boolean, throw, implements, import, this},
  ndkeywordstyle=\color{darkgray}\bfseries,
  identifierstyle=\color{black},
  sensitive=false,
  comment=[l]{//},
  morecomment=[s]{/*}{*/},
  commentstyle=\color{purple}\ttfamily,
  stringstyle=\color{red}\ttfamily,
  morestring=[b]',
  morestring=[b]"
}


\makeatletter
\newcommand\footnoteref[1]{\protected@xdef\@thefnmark{\ref{#1}}\@footnotemark}
\makeatother

\begin{document}

\begin{center}
\vspace*{-2cm}
\hfill\includegraphics[width=4cm]{Bilder/logo_FHE}\\[1cm]
{\Huge \Titel}\\[2cm]
{\Huge\scshape \Was}\\[2cm]
% {\large f�r die Pr�fung zum}\\[0.5cm]
% {\Large \Abschluss}\\[0.5cm]
% {\large des Studienganges \Studiengang}\\[0.5cm]
{\large \Studiengang}\\[0.5cm]
{\large an der}\\[0.5cm]
{\large Fachhochschule Erfurt}\\[0.5cm]
{\large von}\\[0.5cm]
{\large\bfseries \Autor}\\[1cm]
{\large Abgabedatum \AbgabeDatum}
\vfill
\end{center}
\begin{tabular}{l@{\hspace{1cm}}l}
Bearbeitungszeitraum             & \Dauer                       \\
Matrikelnummer                   & \MatrikelNummer              \\
% Kurs                             & \Kursbezeichnung             \\
% Ausbildungsfirma                 & \FirmenName                  \\
%                                  & \FirmenStadt                 \\
Betreuer der Masterarbeit    & \BetreuerFirma               \\
Zweitbetreuer der Masterarbeit     & \BetreuerDHBW                \\
\end{tabular}

\newpage
%Seitenumbruch
%%%%%%%%%%%%%%%%%%%%%%%%%%%%%%%%%%%%%%%%%%%%%%%%%%%%%%%%%%%%%%%%%%%%%%%%%%%%%%
%% Descr:       Vorlage für Berichte der DHBW-Karlsruhe, Erklärung
%% Author:      Prof. Dr. Jürgen Vollmer, vollmer@dhbw-karlsruhe.de
%% $Id: erklaerung.tex,v 1.2 2010/07/22 13:30:27 vollmer Exp $
%%%%%%%%%%%%%%%%%%%%%%%%%%%%%%%%%%%%%%%%%%%%%%%%%%%%%%%%%%%%%%%%%%%%%%%%%%%%%%%

% In Bachelorarbeiten muss eine schriftliche Erklärung abgegeben werden. In allen anderen
% Arbeiten entf�llt diese. Hierin best�tigen die Studierenden, dass die Bachelorarbeit
% selbst�ndig verfasst und s�mtliche Quellen und Hilfsmittel angegeben sind. Diese Erkl�rung
% bildet das zweite Blatt der Arbeit. Der Text dieser Erkl�rung muss auf einer separaten Seite
% wie unten angegeben lauten.

\newpage
\thispagestyle{empty}
\begin{framed}
\begin{center}
\Large\bfseries Erkl\"arung
\end{center}

\noindent
Ich, \Autor, versichere hiermit, dass ich die vorliegende Masterarbeit mit dem
Thema
"`\Titel"'
selbstst�ndig und nur unter Verwendung der angegebenen Quellen und Hilfsmittel angefertigt
habe.

\vspace{3cm}
\noindent
\underline{\hspace{4cm}}\hfill\underline{\hspace{6cm}}\\
Ort~~~~~Datum\hfill Unterschrift\hspace{4cm}
\end{framed}

%%%%%%%%%%%%%%%%%%%%%%%%%%%%%%%%%%%%%%%%%%%%%%%%%%%%%%%%%%%%%%%%%%%%%%%%%%%%%%%
\endinput
%%%%%%%%%%%%%%%%%%%%%%%%%%%%%%%%%%%%%%%%%%%%%%%%%%%%%%%%%%%%%%%%%%%%%%%%%%%%%%%

\newpage
\begin{spacing}{0.9}

%Einf�gen Inhaltsverzeichnis
\tableofcontents
\newpage
\section{Abstract}
\newpage
\section{Einleitung}
In den letzten zwanzig Jahren hat sich die Bedienung von Computern grundlegend ge�ndert. Vor nicht all zu vielen Jahren gab es nur die M�glichkeit mit Hilfe von Maus und Tastatur mit einem Computer zu interagieren. 

Mit dem Aufkommen von Touch-Screens jedoch �nderte sich auch die Benutzung von Computern. Es war nun m�glich, direkt mit dem Bildschirm zu interagieren ohne den Umweg �ber die Maus.

Als dann wenig sp�ter die ersten Sprachsteuerungen auf den Markt kamen, �nderte sich die Interaktion mit dem Computer erneut. So ist es nun m�glich, Computern mittels Sprache Befehle zu erteilen oder Texte zu sprechen, welche automatisch transkribiert werden.

Heute ist es mit manchen Smartphones schon m�glich, Nachrichten wie SMS zu schreiben und zu versenden, ohne das Telefon �berhaupt in die Hand zu nehmen. Dies ist nur durch neuste Entwicklungen in der \ac{KI} m�glich.

Im selben Zeitraum hat sich parallel auch die \ac{VR} entwickelt. Sie versucht Menschen mit Hilfe von verschiedenen Brillen in eine virtuelle Welt zu versetzen. Die Entwicklungen solcher \ac{VR} Systeme wurde vor allem von der Spieleindustrie getrieben, da sie immer neue Versuche unternimmt, Spieler besser in die Welt des Spiels zu versetzen. 

Eine Abstufung der \ac{VR} ist die \ac{AR}, welche versucht die analoge, reale Welt um digitale Inhalte zu erweitern. Hierbei sind die M�glichkeiten f�r den Einsatz von \ac{AR} fast unbegrenzt. 
Es ist also quasi m�glich, jede menschliche T�tigkeit durch die \ac{AR} zu unterst�tzen. 

Genau hier setzt der Schwerpunkt der Arbeit an. Im Verlauf soll versucht werden, ein \ac{NUI} unter Zuhilfenahme moderner \ac{AR} und \ac{KI} Technologien zu erstellen.

Hierf�r werden aktuelle \ac{AR} und \ac{KI} Systeme daraufhin untersucht, wie sie im Zusammenspiel ein \ac{NUI} bilden k�nnen, welches allein durch Sprache und Gesten mit einem Computer interagiert.

Es soll eine �bliche T�tigkeit mit Hilfe einer \ac{AR} Technologie in der realen Welt abgebildet werden, welche dann von einem Menschen nur unter Zuhilfenahme von Sprache und Gesten gesteuert und bearbeitet werden kann.
\newpage
\section{Virtual Reality}
\section{Artificial Intelligence}
\section{Beschreibung eines Natural User Interface}
\section{Versuch der Entwicklung eines Natural User Interface anhand von ... und ...}
\section{Fazit}
\section{zusammenfassung}


\end{spacing}

\newpage
\section{Abk�rzungsverzeichnis}
\begin{acronym}
  \acro{KI}{\emph{k�nstliche Intelligenz}}
  \acro{VR}{\emph{Virtuelle Realit�t}}
  \acro{AR}{\emph{Augumented Reality}}
  \acro{NUI}{\emph{Natural User Interface}}
  \acro{AV}{\emph{Augmented Virtuality}}
  \acro{SDK}{\emph{Software Development Kit}}
  \acro{AI}{\emph{Artificial Intelligence}}
  \acro{AVS}{\emph{Amazon Voice Service}}
  \acro{SRGS}{\emph{Speech Recognition Grammar Specification}}
  \acro{W3C}{\emph{World Wide Web Consortium}}
  \acro{XML}{\emph{Extensible Markup Language}}
  \acro{UWP}{\emph{Universal Windows Plattform}}
\end{acronym}
\newpage
\listoffigures
\newpage
\listoftables
% Abk�rzungsverzeichnis
\newpage
\bibliographystyle{alpha}
% verzeichnis im DIN format
\bibliography{Quellen}
\end{document}
